\begin{proposition}[Upper Bounds Depend on the Ambient Order]
\label{prop:upper-bounds-ambient-order}
Let $(S,\le_S)$ and $(T,\le_T)$ be ordered sets, and let $A\subseteq S\cap T$. Whether $u$ is an upper bound of $A$ must be interpreted relative to the chosen ambient order.
\hyperref[prf:upper-bounds-ambient-order]{Go to proof.}
\end{proposition}
\end{propositionbox}

\begin{remark*}[Standard quantified statement]
\[
\begin{aligned}
&u\text{ is an upper bound of }A\text{ in }S
\Longleftrightarrow
u\in S\text{ and }\forall x\in A\;(x\le_S u),\\
&u\text{ is an upper bound of }A\text{ in }T
\Longleftrightarrow
u\in T\text{ and }\forall x\in A\;(x\le_T u).
\end{aligned}
\]
\end{remark*}

\begin{remark*}[Predicate reading]
\[
\begin{aligned}
&P_S=\mathsf{OrderedSet}(S,\le_S),\qquad
P_T=\mathsf{OrderedSet}(T,\le_T),\qquad A\subseteq S\cap T,\\
&\operatorname{UpperBound}(u,A,P_S)
\Longleftrightarrow
u\in S\land \forall x\in A\;(x\le_S u),\\
&\operatorname{UpperBound}(u,A,P_T)
\Longleftrightarrow
u\in T\land \forall x\in A\;(x\le_T u).
\end{aligned}
\]
\end{remark*}

\begin{remark*}[Negated quantified statement]
\[
\text{The phrase ``}u\text{ is an upper bound of }A\text{''}
\]
is not well-determined until the ambient ordered set and its order relation are fixed.
\end{remark*}

\begin{remark*}[Negation predicate reading]
\[
\begin{aligned}
&P_S=\mathsf{OrderedSet}(S,\le_S),\qquad
P_T=\mathsf{OrderedSet}(T,\le_T),\\
&X=\text{``$u$ is an upper bound of $A$''},\\
&\neg\operatorname{WellDefined}(X,\text{ambient omitted})
\end{aligned}
\]
when the ambient order is omitted.
\end{remark*}

\begin{remark*}[Failure modes]
\begin{description}
  \item[Exposition.]
Ambiguity occurs when the same set $A$ is considered inside different ordered ambient sets or under different order relations. The candidate bound may belong to one ambient set and not another, or the relevant comparison relation may change.


  \item[\operatorname{WellDefined}.]
  \[
\text{The phrase ``}u\text{ is an upper bound of }A\text{''}
\]
\[
\begin{aligned}
&P_S=\mathsf{OrderedSet}(S,\le_S),\qquad
P_T=\mathsf{OrderedSet}(T,\le_T),\\
&X=\text{``$u$ is an upper bound of $A$''},\\
&\neg\operatorname{WellDefined}(X,\text{ambient omitted})
\end{aligned}
\]
\end{description}
\end{remark*}

\begin{remark*}[Failure modes]
\begin{description}
  \item[Exposition.]
\[
\operatorname{UpperBound}_S(u,A)
\text{ and }
\operatorname{UpperBound}_T(u,A)
\]
are separate assertions unless $S=T$ and $\le_S=\le_T$ on the elements being compared.


  \item[\operatorname{WellDefined}.]
  \[
\text{The phrase ``}u\text{ is an upper bound of }A\text{''}
\]
\[
\begin{aligned}
&P_S=\mathsf{OrderedSet}(S,\le_S),\qquad
P_T=\mathsf{OrderedSet}(T,\le_T),\\
&X=\text{``$u$ is an upper bound of $A$''},\\
&\neg\operatorname{WellDefined}(X,\text{ambient omitted})
\end{aligned}
\]
\end{description}
\end{remark*}

\begin{remark*}[Interpretation]
An upper bound is not only about the subset; it is also about where the subset lives. Changing the ambient order changes the universe of possible bounds and the comparison used to test them.
\end{remark*}

\begin{dependencies}
\begin{itemize}
  \item \hyperref[def:ordered-set]{Ordered Set}
  \item \hyperref[def:subset]{Subset}
  \item \hyperref[def:real-upper-bound]{Upper Bound}
\end{itemize}
\end{dependencies}


